\documentclass[10pt,a4paper]{amsart}
\usepackage{amsmath,amssymb,amsthm,graphicx}
\usepackage{verbatim}
\newenvironment{reference}{\comment}{\endcomment}
%\newenvironment{reference}{}{}
\newenvironment{slogan}{\comment}{\endcomment}
\newenvironment{history}{\comment}{\endcomment}

% For commutative diagrams we use Xy-pic
\usepackage[all]{xy}

% We use 2cell for 2-commutative diagrams.
\xyoption{2cell}
% Theorem environments.
%
\theoremstyle{plain}
\newtheorem{theorem}[subsection]{Theorem}
\newtheorem{proposition}[subsection]{Proposition}
\newtheorem{lemma}[subsection]{Lemma}

\theoremstyle{definition}
\newtheorem{definition}[subsection]{Definition}
\newtheorem{example}[subsection]{Example}
\newtheorem{exercise}[subsection]{Exercise}
\newtheorem{situation}[subsection]{Situation}

\theoremstyle{remark}
\newtheorem{remark}[subsection]{Remark}
\newtheorem{remarks}[subsection]{Remarks}

\numberwithin{equation}{subsection}

% Macros
%
\def\lim{\mathop{\mathrm{lim}}\nolimits}
\def\colim{\mathop{\mathrm{colim}}\nolimits}
\def\Spec{\mathop{\mathrm{Spec}}}
\def\Hom{\mathop{\mathrm{Hom}}\nolimits}
\def\Ext{\mathop{\mathrm{Ext}}\nolimits}
\def\SheafHom{\mathop{\mathcal{H}\!\mathit{om}}\nolimits}
\def\SheafExt{\mathop{\mathcal{E}\!\mathit{xt}}\nolimits}
\def\Sch{\mathit{Sch}}
\def\Mor{\mathop{\mathrm{Mor}}\nolimits}
\def\Ob{\mathop{\mathrm{Ob}}\nolimits}
\def\Sh{\mathop{\mathit{Sh}}\nolimits}
\def\NL{\mathop{N\!L}\nolimits}
\def\CH{\mathop{\mathrm{CH}}\nolimits}
\def\proetale{{pro\text{-}\acute{e}tale}}
\def\etale{{\acute{e}tale}}
\def\QCoh{\mathit{QCoh}}
\def\Ker{\mathop{\mathrm{Ker}}}
\def\Im{\mathop{\mathrm{Im}}}
\def\Coker{\mathop{\mathrm{Coker}}}
\def\Coim{\mathop{\mathrm{Coim}}}

% Boxtimes
%
\DeclareMathSymbol{\boxtimes}{\mathbin}{AMSa}{"02}

%
% Macros for moduli stacks/spaces
%
\def\QCohstack{\mathcal{QC}\!\mathit{oh}}
\def\Cohstack{\mathcal{C}\!\mathit{oh}}
\def\Spacesstack{\mathcal{S}\!\mathit{paces}}
\def\Quotfunctor{\mathrm{Quot}}
\def\Hilbfunctor{\mathrm{Hilb}}
\def\Curvesstack{\mathcal{C}\!\mathit{urves}}
\def\Polarizedstack{\mathcal{P}\!\mathit{olarized}}
\def\Complexesstack{\mathcal{C}\!\mathit{omplexes}}
% \Pic is the operator that assigns to X its picard group, usage \Pic(X)
% \Picardstack_{X/B} denotes the Picard stack of X over B
% \Picardfunctor_{X/B} denotes the Picard functor of X over B
\def\Pic{\mathop{\mathrm{Pic}}\nolimits}
\def\Picardstack{\mathcal{P}\!\mathit{ic}}
\def\Picardfunctor{\mathrm{Pic}}
\def\Deformationcategory{\mathcal{D}\!\mathit{ef}}
\usepackage{cleveref}
\crefname{originalreference}{original reference}{original references}
\Crefname{originalreference}{Original reference}{Original references}
\makeatletter
\newlabel{spaces-lemma-category-of-spaces-over-smaller-base-scheme}{{04SG}{source}{}{sourcepointer-1}{}}
\newlabel{spaces-lemma-category-of-spaces-over-smaller-base-scheme@cref}{{[originalreference][0][]04SG}{[1][0][]source}{}{}{}}
\makeatother
\crefname{originalreference}{original reference}{original references}
\Crefname{originalreference}{Original reference}{Original references}
\makeatletter
\newlabel{spaces-lemma-descent-representable-transformations-property}{{03KD}{source}{}{sourcepointer-2}{}}
\newlabel{spaces-lemma-descent-representable-transformations-property@cref}{{[originalreference][0][]03KD}{[1][0][]source}{}{}{}}
\makeatother
\crefname{originalreference}{original reference}{original references}
\Crefname{originalreference}{Original reference}{Original references}
\makeatletter
\newlabel{schemes-lemma-affine-covering-fibre-product}{{01JS}{source}{}{sourcepointer-3}{}}
\newlabel{schemes-lemma-affine-covering-fibre-product@cref}{{[originalreference][0][]01JS}{[1][0][]source}{}{}{}}
\makeatother
\crefname{originalreference}{original reference}{original references}
\Crefname{originalreference}{Original reference}{Original references}
\makeatletter
\newlabel{schemes-lemma-quasi-compact-affine}{{01K4}{source}{}{sourcepointer-4}{}}
\newlabel{schemes-lemma-quasi-compact-affine@cref}{{[originalreference][0][]01K4}{[1][0][]source}{}{}{}}
\makeatother
\title{Characterizations of quasi-separated algebraic spaces by quasi-compactness of fibre products}
\author{The Stacks Project Authors}
\date{stacks-project:spaces-properties.tex:local-snapshot-20261012; GNU FDL 1.2 or later}
\begin{document}
\maketitle
\noindent\textit{Editorial note.} All delivered mathematical text is the author's own, copied verbatim from the pinned Stacks Project chapter file spaces-properties.tex; no statement or proof step was rewritten, reordered or removed.\par\medskip
\noindent\textit{Editorial note.} Span recomputation by label pairing (LANE\_BRIEF 10.2): in this chapter the author's proof environment follows the statement environment as a standalone block; the spec covers the complete environment pair (statement span, proof span) as recorded in the pool candidate.\par\medskip
\noindent\textit{Editorial note.} Occupation scan: combined claim\_key + statement\_anchor + generic-line scan over all registration files (active and withdrawn) found 13 S10-geom published cards for spaces-properties.tex (prop-occupied.json, 2026-10-14); those 13 labels are excluded from this batch by design. For the 77 free labels: 0 occupied. The registry mentions field (7557 files) and focused-seeds (645) both point to the same S10-geom cards (dual-source check), no other occupation.\par\medskip
\ifdefined\hypertarget\hypertarget{sourcepointer-1}{}\fi
\noindent\textit{Source reference.} Original reference 04SG: sources/stacks-chow/tags.txt, source lines 6219--6219. Cross-chapter standard prerequisite cited by the author inside the counted span; it is used as a pointer and its own proof is not part of this item.\par\medskip
\ifdefined\hypertarget\hypertarget{sourcepointer-2}{}\fi
\noindent\textit{Source reference.} Original reference 03KD: sources/stacks-chow/tags.txt, source lines 4673--4673. Cross-chapter standard prerequisite cited by the author inside the counted span; it is used as a pointer and its own proof is not part of this item.\par\medskip
\ifdefined\hypertarget\hypertarget{sourcepointer-3}{}\fi
\noindent\textit{Source reference.} Original reference 01JS: sources/stacks-chow/tags.txt, source lines 2025--2025. Cross-chapter standard prerequisite cited by the author inside the counted span; it is used as a pointer and its own proof is not part of this item.\par\medskip
\ifdefined\hypertarget\hypertarget{sourcepointer-4}{}\fi
\noindent\textit{Source reference.} Original reference 01K4: sources/stacks-chow/tags.txt, source lines 2037--2037. Cross-chapter standard prerequisite cited by the author inside the counted span; it is used as a pointer and its own proof is not part of this item.\par\medskip
\begin{lemma}
\label{lemma-characterize-quasi-separated}
Let $S$ be a scheme. Let $X$ be an algebraic space over $S$.
The following are equivalent
\begin{enumerate}
\item $X$ is a quasi-separated algebraic space,
\item for $U \to X$, $V \to X$ with $U$, $V$ quasi-compact schemes
the fibre product $U \times_X V$ is quasi-compact,
\item for $U \to X$, $V \to X$ with $U$, $V$ affine
the fibre product $U \times_X V$ is quasi-compact.
\end{enumerate}
\end{lemma}

\begin{proof}
Using Spaces, Lemma
\ref{spaces-lemma-category-of-spaces-over-smaller-base-scheme}
we see that we may assume $S = \Spec(\mathbf{Z})$.
Since $U \times_X V = X \times_{X \times X} (U \times V)$
and since $U \times V$ is quasi-compact if $U$ and $V$ are so,
we see that (1) implies (2). It is clear that (2) implies (3).
Assume (3). Choose a scheme $W$ and a surjective \'etale morphism
$W \to X$. Then $W \times W \to X \times X$ is surjective \'etale.
Hence it suffices to show that
$$
j : W \times_X W = X \times_{(X \times X)} (W \times W) \to W \times W
$$
is quasi-compact, see Spaces, Lemma
\ref{spaces-lemma-descent-representable-transformations-property}.
If $U \subset W$ and $V \subset W$ are affine opens, then
$j^{-1}(U \times V) = U \times_X V$ is quasi-compact by assumption.
Since the affine opens $U \times V$ form an affine open covering of
$W \times W$ (Schemes, Lemma \ref{schemes-lemma-affine-covering-fibre-product})
we conclude by
Schemes, Lemma \ref{schemes-lemma-quasi-compact-affine}.
\end{proof}

\end{document}
